\subsection{向量空间的初始拓扑}

命题 7. --- 设 \((\mathrm{E}_{\mathfrak{t}})_{\mathfrak{t}\in\mathbf{I}}\) 是拓扑除环 K 上的一族拓扑向量空间。设 E 是 K 上的向量空间，且对每个 \(\mathfrak{t}\in I\)，设 \(f_{t}\) 是 E 到 \(E_{t}\) 中的线性映射。则 E 上使每个函数 \(f_{t}\) 连续的最粗拓扑是 E 上与 E 的向量空间结构相容的拓扑 T。此外，若对每个 \(x\in E\)，\(\phi(x)\) 表示点 \((f_{\mathfrak{t}}(x))\)（乘积空间 \(F=\prod_{\mathfrak{t}\in I}E_{\mathfrak{t}}\) 中的点），则拓扑 T 是 F 的子空间 \(\phi(\mathrm{E})\) 的拓扑在线性映射 \(\phi\) 下的逆像。

命题的最后一部分是 GT, I, § 4.1, prop. 3 的特例。于是命题由下面的引理得出。

引理. --- 设 M 与 N 是两个向量空间，g 是 M 到 N 中的线性映射。若 \(T_{0}\) 是与 N 的向量空间结构相容的拓扑，则 \(T_{0}\) 在 g 下的逆像与 M 的向量空间结构相容。

例如我们证明，\((\lambda, x) \mapsto \lambda x\) 在每一点 \((\lambda_{0}, x_{0})\)（\(K \times M\) 中的点）处连续。令令 \(y_{0} = g(x_{0})\)。M 中 0 的每个邻域都包含形如 \(g(U)\) 的邻域，其中 U 是 N 中 0 的邻域；由假设，存在 K 中 0 的一个邻域 V 及 N 中 0 的一个邻域 W，使得关系式 \(\lambda - \lambda_{0} \in V\) 与 \(y - y_{0} \in W\) 蕴含 \(\lambda y - \lambda_{0}y_{0} \in U\)。于是关系式 \(\lambda - \lambda_{0} \in V\)、\(x - x_{0} \in g(W)\) 蕴含 \(\lambda x - \lambda_{0}x_{0} \in g(U)\)。类似地可以证明 \((x, y) \mapsto x - y\) 在 \(M \times M\) 中连续。

对每个指标 \(\iota \in I\)，设 \(\mathfrak{B}_{\iota}\) 是 \(E_{\iota}\) 中 0 的一个基本邻域系。由拓扑 \(\mathcal{T}\) 的定义，该拓扑下 0 的邻域滤子由诸族 \(\overline{f}_{\iota}^{1}(\mathfrak{B}_{\iota})\) 中的集合之并生成；换言之，形如 \(\bigcap_{k} \overline{f}_{\iota_k}^{-1}(V_{\iota_k})\) 的集合构成 \(\mathcal{T}\) 下 0 的一个基本邻域系，其中 \((\iota_k)_{1 \leqslant k \leqslant n}\) 是 I 的任意有限指标序列，而对每个指标 \(k\)，\(V_{\iota_k}\) 是 \(\mathfrak{B}_{\iota_k}\) 中的任意集合。

推论 1. --- 设 G 是 K 上的拓扑向量空间。为使 G 到 E 中的映射所成的集 H 等度连续，必须且只需：对所有 \(\iota \in I\)，当 u 在 H 中变化时 \(f_{\iota} \circ u\) 所成的集等度连续。

这是 GT, X, § 2.2, prop. 3 的特例。

推论 2. --- 若诸空间 \(E_{t}\) 是 Hausdorff 的，则为使 \(\mathcal{T}\) 是 Hausdorff 的，必须且只需：对每个 \(x \neq 0\)（在 \(E\) 中），存在指标 \(t \in I\) 使 \(f_{t}(x) \neq 0\)。

因为这时 \(\phi(E)\) 是 Hausdorff 空间，而为使 \(\mathcal{T}\) 是 Hausdorff 的，显然必须且只需 \(\phi\) 是单射；注意这时可以把 \(E\)（带 \(\mathcal{T}\)）与子空间 \(\phi(E)\)（\(\prod_{i} E_{i}\) 的子空间）借助于映射 \(\phi\) 等同起来。

推论 3. --- 设诸 \(E_{t}\) 完备且 \(\phi(E)\) 在 \(F = \prod_{t \in I} E_{t}\) 中闭。则 E 在拓扑 T 下完备。

因为这时子空间 \(\phi(E)\)（\(F\) 的子空间）完备 (GT, II, § 3.4, prop. 8 与 § 3.5, prop. 10)，因此在逆像拓扑下 \(E\) 亦然 (GT, I, § 7.6, prop. 10 及 GT, II, § 3.1, prop. 4)。

* 例. --- 设 \(\mathcal{D}'(\mathbf{R})\) 是 \(\mathbf{R}\) 上的广义函数空间；设 \(p\) 是满足 \(1 \leqslant p \leqslant +\infty\) 的数，令 \(j: \mathrm{L}^p(\mathbf{R}) \to \mathcal{D}'(\mathbf{R})\) 为典范单射，它是连续的（当 \(\mathrm{L}^p(\mathbf{R})\) 赋以其赋范空间拓扑而 \(\mathcal{D}'(\mathbf{R})\) 赋以强拓扑时）。对每个广义函数 \(f \in \mathcal{D}'(\mathbf{R})\)，把 \(f\) 的导数记作 \(\mathrm{D}(f)\)；回忆 \(f \mapsto \mathrm{D}(f)\) 是 \(\mathcal{D}'(\mathbf{R})\) 的连续自同态。于是设 E 是 \(\mathrm{L}^p(\mathbf{R})\) 的向量子空间，由那些 \(f \in \mathrm{L}^p(\mathbf{R})\)（满足 \(\mathrm{D}(f) \in \mathrm{L}^p(\mathbf{R})\)）组成，并在 E 上赋予使典范单射 \(i: \mathrm{E} \to \mathrm{L}^p(\mathbf{R})\) 与映射 \(\mathrm{D}: \mathrm{E} \to \mathrm{L}^p(\mathbf{R})\) 都连续的最粗拓扑（\(\mathrm{L}^p(\mathbf{R})\) 赋以其赋范空间拓扑）。对这个拓扑，空间 E 是完备的。因为，E 在 \(\mathrm{F} = \mathrm{L}^p(\mathbf{R}) \times \mathrm{L}^p(\mathbf{R})\) 中经映射 \(\phi: f \mapsto (f, \mathrm{D}(f))\) 所成的像是闭的，因为它是像 G 在 \(\mathrm{L}^p(\mathbf{R}) \times \mathrm{L}^p(\mathbf{R})\) 上的迹，其中 G 是 \(\mathcal{D}'(\mathbf{R})\) 在 \(\mathcal{D}'(\mathbf{R}) \times \mathcal{D}'(\mathbf{R})\) 中经映射

\[
\phi_ {0}: f \mapsto (f, \mathrm{D} (f));
\]

所成的像。现在 \(\mathbf{G}\) 是 \(\phi_0\) 的图像，故在 \(\mathcal{D}'(\mathbf{R}) \times \mathcal{D}'(\mathbf{R})\) 中是闭的（GT, I, § 8.1, prop. 2 的 cor. 2）；而 \(\phi(\mathbf{E})\) 是 \(\mathbf{G}\) 经 \(i \times i\)（它连续）的原像，于是我们看出 \(\phi(\mathbf{E})\) 在 F 中是闭的。*

推论 4. --- 设 E 是拓扑除环 K 上的一个向量空间，并设 \((\mathcal{T}_{\iota})_{\iota\in I}\) 是与 E 的向量空间结构相容的一族拓扑；则诸拓扑 \(T_{\iota}\) 的上确界 T 与 E 的向量空间结构相容。

因为，若用 \(E_{t}\) 表示由拓扑 \(T_{t}\) 从 E 得到的拓扑向量空间，并用 \(f_{t}\) 表示 E 到 \(E_{t}\) 上的恒等映射，则 T 是使诸 \(f_{t}\) 连续的最粗拓扑。

